% General packages: 
\documentclass{article}
\usepackage[utf8]{inputenc}
\usepackage{graphicx}
\usepackage{amsmath}
\usepackage{indentfirst}
\usepackage{amsthm}
\usepackage{fourier}
% Algorithms:
\usepackage[algo2e,boxruled, linesnumbered]{algorithm2e}
\SetKwComment{Comment}{//}{}
\SetNlSty{}{}{:}



\begin{document}

\begin{algorithm2e}[H]
\caption{Powersort(A[1...n])}\label{ps1}
$X := $ stack of runs\Comment*[r]{{capacity$\lfloor{lgn(n)}\rfloor + 1$}}
$P :=$ stack of powers\Comment*[r]{capacity $\lfloor{lgn(n)}\rfloor + 1$}
$s_1 := 1; e_1 = ExtendRunRight(A[1], n)$\;
\While{$e_1 < n$}{
     $s_2 := e_1 + 1;$ $e_2 := ExtendRunRight(A[s_2], n)$ \Comment*[r]{A[$s_2..e_2$] next run}
     $p := NodePower(s1, e1, s2, e2, n)$\Comment*[r]{{\small $P_j$ for node j between A[$s_1..e_1$] and A[$s_2..e_2$]}}
     \While {P.top() $>$ p}{
        \nlset{}\Comment{previous merge deeper in tree than current}
        P.pop()\Comment*[r]{$\rightarrow$ merge and replace run A[$s_1..e_1$] by result}
        ($s_1, e_1$) := Merge(X.pop(), A[$s_1..e_1$]);
     }
     X.push(A[s1, e1])\; 
     P.push(p)\;
     $s_1 := s_2; e_1 := e_2$\Comment*{Now A[$s_1..s_1$] is the rightmost run}}
 \While{¬X.empty()}{
     ($s_1, e_1$) := Merge(X.pop(), A[$s_1..e_1$]);
 }
\end{algorithm2e}

\end{document}
% \newpage
% \begin{algorithm}
% \caption{First\_{}C-Bounded\_{}Powersort(A[1...n])}\label{ps2}
% \begin{algorithmic}
% $X :=$ C-bounded stack of runs \Comment*[r]{capacity C}
% $P :=$ C-bounded stack of powers \Comment*[r]{capacity C}
% $s_1 := 1; e_1 = ExtendRunRight(A[1], n)$
% \While{$e_1 < n$} \Comment*[r]{{\small This while can be tweaked to use max amount of passes}}
%     \While{$e_1 < n$} \Comment*[r]{We do a pass covering the entire array}
%         $s_2 := e_1 + 1; e_2 := ExtendRunRight(A[s_2], n)$
%         $p := NodePower(s1, e1, s2, e2, n)$ \Comment*[r]{{\small Using our node power algorithm}}
%         \While {P.top() $>$ p} \Comment*[r]{{\small We follow the same merge policy as Munro and Wilde}}
%             P.pop()
%             ($s_1, e_1$) := Merge(X.pop(), A[$s_1..e_1$])
%         \EndWhile 
%         X.push(A[s1, e1]); P.push(p)
%         $s_1 := s_2; e_1 := e_2$
%     \EndWhile 
%     \While{¬X.empty()}
%         ($s_1, e_1$) := Merge(X.pop(), A[$s_1..e_1$]) 
%     \EndWhile \Comment*[r]{{\small We're using a limited stack, so we need to check for another pass}}
%     $s_1 := 1; e_1 = ExtendRunRight(A[1], n)$ \Comment*[r]{{\small If $e_1$ = end, the array is sorted}}
% \EndWhile 
% \end{algorithmic}
% \end{algorithm}

% \newpage
% \begin{algorithm}
% \caption{C-Bounded\_{}Powersort(A[1...n])}\label{ps3}
% \begin{algorithmic}
% $X :=$ C-bounded stack of runs and powers
% $s_1 := 1; e_1 = ExtendRunRight(1, n)$
% \While{$true$}
%     \If{$e_1 == n$}
%     break
%     \EndIf
%     \While{$e_1 < n$} 
%         $s_2 := e_1 + 1; e_2 := ExtendRunRight(s_2, n)$ 
%         $p := NodePower(s1, e1, s2, e2, n)$
%         \While {X.size \textgreater \text{ }1 and X.top().power $>$ p}
%             ($s_1, e_1$) := Merge(X.top().boundary, A[$s_1..e_1$])
%             X.pop()
%         \EndWhile X.push(s1, p);
%         $s_1 := s_2; e_1 := e_2$
%     \EndWhile
%     \While{$X.size > 1$}
%         ($s_1, e_1$) := Merge(X.pop(), A[$s_1..e_1$]) 
%     \EndWhile
%     $s_1 := n; e_1 = ExtendRunLeft(A[n], 1)$ \Comment*[r]{{\scriptsize Now, we're scanning from right to left!}}
%     \If{$e_1 == 1$}
%         break \Comment*[r]{{\scriptsize If we could extend until the beginning, our run is the entire array.}}
%     \EndIf
%     \While{$e_1 > 1$} \Comment*[r]{We're scanning from n down to 1}
%         $s_2 := e_1 - 1; e_2 := ExtendRunLeft(s_2, 1)$ 
%         $p := NodePower(s1, e1, s2, e2, n)$
%         \While {X.size \textgreater \text{ }1 and X.top() $>$ p}
%             ($s_1, e_1$) := Merge(X.top().boundary, A[$s_1..e_1$])
%             X.pop()
%         \EndWhile 
%         X.push(s1, p); 
%         $s_1 := s_2; e_1 := e_2$
%     \EndWhile 
%     \While{$X.size > 1$}
%         ($s_1, e_1$) := Merge(X.pop(), A[$s_1..e_1$]) 
%     \EndWhile
%     $s_1 := 1; e_1 = ExtendRunRight(1, n)$ \Comment*[r]{{\small Now, we're scanning from left to right!}}
% \EndWhile
% \end{algorithmic}
% \end{algorithm}
